% year: תשפ"ד
% semester: א
% moed: א
% number: 4
% points: 14
% categories: taylor
% source: exams/מבחנים/תשפד/מועד א תשפד.pdf

\begin{question}
העריכו בעזרת קירוב לינארי את $\sqrt[3]{30}$, ותנו חסם לשגיאה.
\end{question}

\begin{hint}
השתמשו ב-$f(x)=\sqrt[3]{x}$ סביב הנקודה $x_0=27$, שבה השורש ידוע.
\end{hint}

\begin{hint}
לחסם השגיאה השתמשו בשארית לגרנז' של פולינום טיילור מסדר 1: $R_1(x)=\frac{f''(c)}{2}(x-x_0)^2$.
\end{hint}

\begin{solution}
נבחר $f(x)=\sqrt[3]{x}=x^{1/3}$ ו-$x_0=27$ (נקודה קרובה ל-$30$ שבה $f(27)=3$).

\textbf{הקירוב הלינארי.} $f'(x)=\frac13x^{-2/3}$, ולכן $f'(27)=\frac{1}{3\cdot 9}=\frac1{27}$. הקירוב הלינארי (פולינום טיילור מסדר 1):
\[ f(x)\approx f(27)+f'(27)(x-27)=3+\frac{x-27}{27}, \]
ובפרט
\[ \sqrt[3]{30}\approx 3+\frac{3}{27}=3+\frac19=\frac{28}{9}\approx 3.1111. \]

\textbf{חסם לשגיאה.} לפי משפט טיילור עם שארית לגרנז', קיימת $c\in(27,30)$ כך ש-
\[ \sqrt[3]{30}-\frac{28}{9} = R_1(30) = \frac{f''(c)}{2!}(30-27)^2. \]
$f''(x)=-\frac29x^{-5/3}$, ולכן
\[ |R_1(30)| = \frac{2}{9}c^{-5/3}\cdot\frac{9}{2} = \frac{1}{c^{5/3}} < \frac{1}{27^{5/3}}=\frac{1}{3^5}=\frac1{243}\approx 0.0041, \]
כי $c>27$ והפונקציה $c^{-5/3}$ יורדת.

בנוסף, $f''<0$ ולכן $R_1<0$: הקירוב הוא הערכת-יתר. לסיכום
\[ \boxed{\sqrt[3]{30}\approx\frac{28}{9}\approx3.111,\qquad 0<\frac{28}{9}-\sqrt[3]{30}<\frac{1}{243}\approx0.0041.} \]
(לבדיקה: $\sqrt[3]{30}\approx3.1072$, והשגיאה בפועל כ-$0.0039$.)
\end{solution}
